\documentclass[11pt,a4paper]{article}
\usepackage[margin=2.5cm]{geometry}
\usepackage[T1]{fontenc}
\usepackage{booktabs}
\usepackage{xcolor}
\usepackage{listings}
\usepackage[hidelinks]{hyperref}

\emergencystretch=2em
\definecolor{codebg}{gray}{0.95}
\lstset{language=C++, basicstyle=\ttfamily\small, keywordstyle=\color{blue!70!black},
  commentstyle=\color{green!50!black}\itshape, backgroundcolor=\color{codebg},
  frame=single, rulecolor=\color{black!20}, breaklines=true, showstringspaces=false,
  tabsize=2, columns=flexible}

\title{ATmega32A Module \\ \Large Seven-Segment Display}
\author{Ali Sahafi}
\date{}

\begin{document}
\maketitle

\section{Theory}

A seven-segment display is simply \textbf{eight LEDs} (segments
\emph{a}--\emph{g} plus the decimal point) arranged as a digit. In a
\textbf{common-anode} display all LED anodes share the VCC pin, so a segment
lights up when its cathode pin is pulled \textbf{LOW} --- the segments are
\emph{active-low}. Displaying ``5'' just means turning on the right subset of
segments (a, f, g, c, d).

This driver is hardwired to \textbf{PORTB}: one call writes the complete
segment pattern for a digit.

\section{Pin mapping}

The driver expects the display wired to PORTB as follows (each segment
through a 330\,$\Omega$ resistor):

\begin{center}
\begin{tabular}{llllllllc}
\toprule
\textbf{PORTB bit} & PB0 & PB1 & PB2 & PB3 & PB4 & PB5 & PB6 & PB7 \\
\midrule
\textbf{Segment}   & e   & d   & c   & b   & f   & g   & a   & dp \\
\bottomrule
\end{tabular}
\end{center}

A bit value of \textbf{0 turns the segment on} (active-low, common-anode).
For example digit 8 is \texttt{0b10000000}: every segment on, decimal
point off.

\section{API reference}

\begin{center}
\begin{tabular}{ll}
\toprule
\textbf{Method} & \textbf{Description} \\
\midrule
\texttt{SevenSeg(digit)} & Display digit 0--9 on PORTB; any other value shows \texttt{E} \\
\bottomrule
\end{tabular}
\end{center}

The function also \emph{returns} the segment pattern it wrote, and PORTB
must be configured as output first:
\lstinline|GPIO.setDirection(DDRB, ALL, OUTPUT);|

\section{Example: counting 0--9}

\paragraph{Wiring.}
\begin{itemize}
  \item Common-anode 7-segment display: common pin to VCC, segments to
        PB0--PB7 via 330\,$\Omega$ resistors (mapping table above).
\end{itemize}

\paragraph{Build \& flash.} \texttt{make sevenseg}

\lstinputlisting{example.cpp}

\section{Exercises}

\begin{enumerate}
  \item Combine this module with GPIO: count one step per \emph{button press}
        instead of every 500\,ms.
  \item Extend \texttt{SevenSeg()} with hexadecimal digits A--F. Derive the
        segment patterns yourself using the pin-mapping table.
  \item What changes if your display is \textbf{common-cathode} instead of
        common-anode? Modify the driver accordingly (hint: one operator).
\end{enumerate}

\end{document}
